\documentclass{amsart}
\usepackage{amsmath}
\usepackage{amssymb}
\usepackage{amsthm}
\usepackage{tikz}

\begin{document}
    \section{Preliminaries}
    \subsection{Decimal Expansion}
    
    \subsubsection{Exercise:} Prove that if the decimal expansion of $r$ is eventually periodic then $r$ is a rational number.
    \begin{proof}
        Assume that the decimal expansion of $r$ \textit{is} eventually periodic with period $p$. By definition, there exists some $m \in \mathbb{N}$ such that

        \[a_{n+p} = a_n \quad \text{for all } n > m\]

        We can express $r$ as follows.

        \[ r = \sum_{i = 1}^{\infty} 10^{-i} a_i\] 

        We can split this sum to isolate the non-periodic prefix:

        \[ r = \sum_{i = 1}^{m} 10^{-i} a_i + \sum_{i = m + 1}^{\infty} 10^{-i} a_i \] 

        Since $S_1 = \sum_{i = 1}^{m} 10^{-i} a_i$ is a finite sum of rationals, $S_1 \in \mathbb{Q}$.

        Now let $S_2 = \sum_{i = m+1}^{\infty} 10^{-i} a_i$. 
        
        \[S_2 = \sum_{i = m+1}^{\infty} 10^{-i} a_i = \sum_{i = 1}^{\infty} 10^{-(i + m)} a_{i + m}\]

        We separate the sum into periodic blocks:
        \begin{align*}
            S_2 &= \sum_{i = 1}^{\infty} 10^{-(i + m)} a_{i + m} \\
            &= \sum_{k = 0}^{\infty} \sum_{i = 1}^{p} 10^{-(i + m + kp)} a_{i + m} \\
            &= \sum_{k = 0}^{\infty} 10^{-pk} \sum_{i = 1}^{p} 10^{-(i + m)} a_{i + m}
        \end{align*}
        
        Let $A = \sum_{i = 1}^{p} 10^{-(i + m)} a_{i + m}$. Since A is a finite sum of rationals, $A \in \mathbb{Q}$, and $S_2 = A \cdot \sum_{k = 0}^{\infty} 10^{-pk}.$ We recognize this as a geometric series with common ratio $10^{-p}$, and since $|10^{-p}| < 1$ for all $p \geq 1$, the series converges to a rational value. Therefore, $S_2$ is a product of rationals, and $r = S_1 + S_2$ is rational.
    \end{proof}
    \newpage
    \section{Baker's transformation and the slope-2 tent map}
    \subsubsection{Exercise}
    \begin{center}[h]
        \begin{tikzpicture}[scale=5]
            % Grid / Axes
            \draw[thin, gray!40, dashed] (0.5, 0) -- (0.5, 1);
            \draw[thin, gray!40, dashed] (0, 1) -- (1, 1);
            \draw[thin, gray!40, dashed] (1, 0) -- (1, 1);
            
            \draw[->] (0, 0) -- (1.15, 0) node[right] {$x$};
            \draw[->] (0, 0) -- (0, 1.15) node[above] {$g(x)$};

            % Axis Ticks
            \node[below] at (0,0) {$0$};
            \node[below] at (0.5,0) {$\frac{1}{2}$};
            \node[below] at (1,0) {$1$};
            \node[left] at (0,1) {$1$};

            % Tent / Baker Plot Lines
            \draw[thick, blue] (0, 0) -- (0.5, 1);
            \draw[thick, blue] (0.5, 0) -- (1, 1);
        \end{tikzpicture}
    \end{center}

    \subsection{Baker’s transformation and binary expansions.}
\end{document}
